\input{../tex-inputs/latex-leseansicht-vorspann}

\section[Gustav Schwarzkopf an Arthur Schnitzler, 21. 6. 1893]{L04275 Gustav Schwarzkopf an Arthur Schnitzler, 21. 6. 1893}
\nopagebreak\mylabel{L04275v}
\rehead{ }\normalsize\beginnumbering\briefempfaengerindex{Schnitzler, Arthur@\textsc{Schnitzler, Arthur}!zzzSchwarzkopf, Gustav@\emph{von Gustav Schwarzkopf}!1893-06-213@{21. 6. 1893}|(be}
\toendnotes[C]{\smallbreak\pagebreak[2]}
\correspDesc{Versand  durch Gustav Schwarzkopf am 21. 6. 1893 in Brühl
\newline{}Erhalt  durch Arthur Schnitzler im Zeitraum [21. 6. 1893
                  – 22. 6. 1893?] in Wien}\toendnotes[C]{\smallbreak}
\Standort{CUL, Schnitzler, B 96a.}
\physDesc{Brief, 1 Blatt, 2 Seiten, 1034 Zeichen
\newline{}Handschrift: schwarze Tinte, deutsche Kurrent}\toendnotes[C]{\smallbreak}
\pstart
           \raggedleft{}{\pb}\textsc{Brühl\oindex{Brühl@\textbf{Brühl}, \emph{Tal}|pw}}{ }21/6 1893\pend
           
\pstart{}Sehr geehrter Herr Doctor!\pend\vspace{0.5em}
\pstart
           Ich beeile mich im Nachſtehenden Ihre \label{K_L04275-1v}\edtext{Fragen}{\lemma{\textnormal{\emph{Fragen}}}\Cendnote{\textnormal{Offenbar hatte Schnitzler
                  um Auskunft gebeten, wo man in Mödling\oindex{Mödling@\textbf{Mödling}, \emph{Hauptstadt}|pwk} am besten übernachten kann. Bereits
                  am Folgetag, dem 22. 6. 1893 zog Schnitzler
                  für eine Nacht ins Hotel Hajek\oindex{Hotel Hajek@\textbf{Hotel Hajek}, \emph{Hotel}|pwk}.}}}\label{K_L04275-1} zu beantworten.\pend
           
\pstart
           In Mödling\oindex{Mödling@\textbf{Mödling}, \emph{Hauptstadt}|pw} kann man am beſten in der »Stadt Mödling\oindex{Hotel Stadt Mödling@\textbf{Hotel Stadt Mödling}, \emph{Hotel}|pw}« oder beim »Lamm\oindex{Deisenhofers Hotel »zum goldenen Lamm«@\textbf{Deisenhofers Hotel »zum goldenen Lamm«}, \emph{Hotel}|pw}« wohnen, doch würde ich dazu nur im äußerſten Fall
               rathen.\pend
           
\pstart
           Bei \textsc{Hajek\oindex{Hotel Hajek@\textbf{Hotel Hajek}, \emph{Hotel}|pw}} bewegen{ }ſich die Preiſe für \uline{1. \textsc{Cabinet}} mit 2. Betten (wenn ein{ }ſolch\textcolor{gray}{es} frei) zwiſchen 18–25 f
               wöchentl. für ein{ }ſolches \introOben{}größeres\introOben{} Zimmer \textsc{ca} 30f. wöchentlich. Bei Schönberger\oindex{Zum goldenen Stern@\textbf{Zum goldenen Stern}, \emph{Lokal}|pwv} iſt bis – 15 Juli gar nichts frei,
               und von da ab nur im Dachcabinet mit 1. Bett.\pend
           
\pstart
           In Erwägung zu ziehen wäre eventuell noch das \textsc{Hotel z. Jordan\oindex{Hotel-Restauration »zum Jordan«@\textbf{Hotel-Restauration »zum Jordan«}, \emph{Hotel}|pw}} auf der Brühlerstraſſe\oindex{Brühler Straße@\textbf{Brühler Straße}, \emph{Meerenge}|pw} oder \substVorne{}\textsuperscript{die}\substDazwischen{}\textsc{Hotel}\substHinten{} zur \textsc{Helmstreit}mühle\oindex{Hotel zur Helmstreitmühle@\textbf{Hotel zur Helmstreitmühle}, \emph{Hotel}|pw}. {\pb}Frau \textsc{Schönberger\pwindex{Schönberger, Christine 14.\,10.\,1848 – 6.\,10.\,1927@\textsc{Schönberger, Christine} (14.\,10.\,1848 – 6.\,10.\,1927), \emph{Gastwirtin}|pw}} hat{ }ſich übrigens bereit erklärt falls Sie{ }ſich ein Privat-Zimmer nehmen (\strikeout{wo} es befinden{ }ſich 2 in nächſter Nähe des Hotels\oindex{Zum goldenen Stern@\textbf{Zum goldenen Stern}, \emph{Lokal}|pwv}) das nötige Bettzeug und
               Wäſche beizuſtellen.\pend
           
\pstart
           Ich würde es demnach wirklich für das beſte halten, wenn Sie Ihr Vorhaben ausführten
               und morgen herauskämen. Es wird dann unſeren vereinten Bemühungen gewiß gelingen
               etwas Paſſendes zu finden.\pend
           
\pstart
           In dieſer Erwartung verbleibe ich mit her\textcolor{gray}{zl}. Grüßen\pend
           
\pstart
           Ihr ganz \textcolor{gray}{errg}{\\[\baselineskip]}\spacefill\mbox{GSchwarzkopf}\pend
           \leftskip=0em{}\selectlanguage{ngerman}\endnumbering\briefempfaengerindex{Schnitzler, Arthur@\textsc{Schnitzler, Arthur}!zzzSchwarzkopf, Gustav@\emph{von Gustav Schwarzkopf}!1893-06-213@{21. 6. 1893}|)be}\mylabel{L04275h}
\begin{anhang}
\end{anhang}\newcommand{\dateiname}{L04275}\newcommand{\titel}{Gustav Schwarzkopf an Arthur Schnitzler, 21. 6. 1893}\newcommand{\editorInnen}{Herausgegeben von Jahnke, SelmaMüller, Martin Anton}\input{../tex-inputs/latex-leseansicht-abspann}
